\begin{document}
\docTitle{Running services}

This page lists the components that are expected to run or be maintained on the MacBook.

\paragraph{Tailscale}
On the MacBook, a userspace \code{tailscaled} daemon is managed by the \path{io.tailscale.tailscaled.userspace} LaunchAgent.
It allows different nodes to communicate through a private VPN.
Additional information is available in the \emph{Tailscale setup} document.

\paragraph{Local Codex session search}
The \code{search-in-the-past} workflow runs locally on the MacBook.
It does not depend on OpenSearch, a remote Linux node, or a tailnet search service.
The raw Codex session files at the following locations remain the source of truth:
\begin{itemize}
    \item \code{/Users/tbe/.codex/sessions}
    \item \code{/Users/tbe/.codex/archived_sessions}
\end{itemize}

The \code{session-search} command maintains a local SQLite FTS5 index stored at:
\begin{DocCode}
~/.codex/local-session-search/session-search.sqlite
\end{DocCode}
The directory and database are restricted to the current macOS user.
The index stores extracted messages and short conversation segments with metadata such as the timestamp, role, current working directory, session ID, absolute source-file path, and source line number.
Conversation, developer, system, and supported tool records are indexed; normal searches query only user and assistant messages.
Tool calls, tool output, commands, errors, and developer or system messages require the explicit \code{tools} or \code{all} search mode.

\paragraph{Index refresh and verification}
Every normal search first performs an incremental refresh.
The refresh scans the two local session roots, skips unchanged \code{rollout-*.jsonl} files, reads only newly appended bytes when possible, fully reindexes rewritten files, and removes derived records for deleted source files.
No background daemon is required.

Common commands are:
\begin{DocCode}
session-search stats
session-search search 'prior decision'
session-search search 'pytest failure' --mode tools
session-search verify 'exact phrase'
session-search rebuild
\end{DocCode}

The raw JSONL files remain authoritative.
When a search result is used to establish what was decided, what was completed, or what remains unfinished, it must be verified against the raw JSONL file and line number.

The \code{verify} command uses ripgrep for exact fixed-string searches.
Its direct equivalent is:
\begin{DocCode}
rg -n -F --glob 'rollout-*.jsonl' '<exact text>' \
  /Users/tbe/.codex/sessions \
  /Users/tbe/.codex/archived_sessions
\end{DocCode}

\paragraph{Out of scope}
This implementation does not require a daemon, Spotlight indexing, OpenSearch, Tantivy, embeddings, or queries to another machine.
The infrastructure consists of a local SQLite FTS5 file and raw \code{rg -F} verification.

\end{document}
